\scanpage{22}

\[
\operatorname{Hom}_k(\Lambda^{\unsure{i|j}}\Psi_{\ast},\Lambda^j\Psi_{\ast})\simeq
\begin{cases}0&\text{si }j\ne i\\k&\text{si }j=i\end{cases}
\]

de \unsure{façon|\deleted{façon}} précise

\[
\operatorname{Hom}_k(\Lambda^i\Psi_{\ast},\Lambda^i\Psi_{\ast})\simeq k\,\mathrm{id}_{\Lambda^i\Psi_{\ast}}\xleftarrow{\sim} k,
\]

et

\[
\operatorname{Hom}_k(\Lambda^i\Psi_{\ast},\Phi_{\deleted{\unsure{j|\bullet}}}(j,n))\simeq
\begin{cases}0&\text{si }j\ne\deleted{\unsure{i,}}i+1\\k&\text{si }j=i+1\end{cases}
\]

en particulier

\[
\operatorname{Hom}_k(\Lambda^i\Psi_{\ast},\Lambda^j\Phi_{\ast})\simeq
\begin{cases}0&\text{si }j\ne i+1\\k&\text{si }j=i+1\end{cases}
\]

et

\[
\operatorname{Hom}_k(\Lambda^i\Psi_{\ast},\Lambda^{i+1}\Phi_{\ast})=k\cdot\unsure{i|\iota|\partial}\xleftarrow{\sim}k
\]

où $i:\Lambda^i\Psi_{\ast}\hookrightarrow\Lambda^{i+1}\Phi_{\ast}$ est l'inclusion canonique.

\uline{NB} On a un \inserted{hom de} suites exactes can

\diagram{\begin{tikzcd}[ampersand replacement=\&,column sep=small,row sep=large]
0\arrow[r,dash] \& \Lambda^{i-1}\Psi_{\ast}\arrow[r]\arrow[d,"n\,\mathrm{id}"'] \& \Lambda^{i\deleted{\unsure{+1}}}\Phi_{\ast}\arrow[r]\arrow[d,equal] \& \Lambda^i\Psi_{\ast}\arrow[r,dash]\arrow[d,Rightarrow,"\mathrm{id}"] \& 0\\
0\arrow[r,dash] \& \Lambda^{i-1}\Psi_{\ast}\arrow[r] \& \deleted{\unsure{\Lambda^i}}\Phi_{\ast}(i,n)\arrow[r,dash] \& \Lambda^i\Psi_{\ast}\arrow[r,dash] \& 0
\end{tikzcd}}

qui montre que $\Phi_{\ast}(i,n)$ est l'ext déduite de \inserted{l'ext} $\Lambda^i\Phi_{\ast}$ de $\Lambda^i\Psi_{\ast}$ par $\Lambda^{i-1}\Psi_{\ast}$ par l'hom $n\,\mathrm{id}:\Lambda^{i-1}\Psi_{\ast}\mathrel{-}\Lambda^{i-1}\Psi_{\ast}$.

% archive sheet number: 21
